\section{Experimental Program}
\label{sec:experiments}
\begin{figure}[t]
\centering
\begin{tikzpicture}[node distance=6mm]
\node[primarybox, text width=0.78\linewidth] (e0) {\textbf{Tier 1: controlled spectral models}\\[-1mm]{\scriptsize Counterexamples, threshold ties, rank deficiency, support overlap, and exact local-rate diagnostics.}};
\node[secondarybox, text width=0.78\linewidth, below=6mm of e0] (e1) {\textbf{Tier 2: controlled neural bridge}\\[-1mm]{\scriptsize Toy superposition networks and small transformers with known latent features.}};
\draw[arrSecondary] (e0) -- (e1);
\node[successbox, text width=0.78\linewidth, below=6mm of e1] (e2) {\textbf{Tier 3: foundation-model validation}\\[-1mm]{\scriptsize Open-model activation subspaces with document-level stability and causal interventions.}};
\draw[arrSecondary] (e1) -- (e2);
\end{tikzpicture}
\caption{Validation stack. Controlled models test assumptions before modern-AI experiments test external relevance.}
\label{fig:validation-stack}
\end{figure}

\begin{table}[t]
\centering
\caption{Pre-specified validation grid.}
\label{tab:a-experiments}
\small
\begin{tabular}{p{0.16\linewidth}p{0.25\linewidth}p{0.23\linewidth}p{0.15\linewidth}}
\toprule
Tier & Manipulated factors & Primary metric & Falsifier \\
\midrule
Synthetic & Margin, gap, overlap, ties, initialization & Identification time and local slope & No relation to assumption diagnostics \\
Neural bridge & Known concepts, superposition, noise & Support/causal recovery & No improvement over matched neurons \\
Foundation model & Layer, domain, model family, support budget & Stability, causal effect, screening cost & Unstable supports or no causal value \\
\bottomrule
\end{tabular}
\end{table}


\subsection{Questions and preregistered tests}
The experiments answer four distinct questions: (Q1) do the corrected objectives differ materially; (Q2) do strict margins predict finite identification; (Q3) do active eigengaps or tangent-Jacobian radii predict local slopes; and (Q4) does screening reduce measured cost without changing the selected limit? Expected patterns are registered before large-model experiments.

\subsection{Tier 1: controlled sparse spectral models}
Use at least 50 independent replicates per condition for small and medium problems. Generate rank-one and multi-spike covariance models with disjoint, overlapping, nested, and component-specific supports. Sweep threshold margin, active eigengap, support overlap, initialization angle, repeated eigenvalues, and the number of components. Include adversarial cases with zero responses, threshold ties, rank-deficient gradient matrices, and nonunique polar factors. For small dimensions, enumerate or solve exact support problems to verify objective values and refits.

Primary plots are: identification iteration versus measured limiting margin; observed contraction versus the relevant spectral prediction; objective and residual trajectories; disagreement between A-SUM and A-SQ; permanent-screening ratio and wall-clock reduction; and raw versus refitted variance. A theory-facing regression should include uncertainty over independent trials and report where assumptions fail.

\subsection{Tier 2: controlled neural bridge}
Train a toy superposition network and a small transformer on concept tasks with known latent factors. Record activations at a fixed site, define the statistically independent unit at the sequence or document level, and test recovery of known channel groups or causal directions. This bridge evaluates whether threshold margins and support overlap remain meaningful after nonlinear training without claiming that the network follows a spiked Gaussian model.

\subsection{Tier 3: transformer activations}
Use at least two open model families and predeclare layers, activation sites, calibration corpora, and document-level splits. Compare dense PCA, deflated and block sparse spectral methods, truncated power, Fantope/convex methods where feasible, gated SAEs, SAEBench-compatible SAE baselines, native-neuron selection, and random supports. All methods receive matched support size and, where meaningful, the same restricted refit.

Report support stability across document and domain splits; reconstruction before and after refit; MIB or causal-variable performance; ablation and activation-patching effects; redundancy across components; correlation evaluations; runtime; and memory. No semantic claim is made from feature labels alone.

\subsection{Ablations and stage-separated estimator lineage}
Every table contains separate rows for raw selection, restricted refit, support-fixed recalibration, and optional end-to-end adaptation. Additional ablations vary initialization, tie policies, polar completion, score-orthogonality versus overlap penalties, screening radius quality, and stopping rules.

\subsection{Falsification}
The central thesis is weakened if corrected and printed objectives agree across every heterogeneous-energy case, measured positive margins fail to predict stabilization, no local spectral phase appears after support fixation, safe screening gives negligible savings, or sparse activation components do not outperform equally sized native-neuron and random-support controls on causal tasks.
